\subsection{LIMITS AND CONTINUITY}

Let X, Y be two topological spaces, f a mapping of X into Y, B the neighbourhood filter in X of a point \(a \in X\) . Instead of saying that \(y \in Y\) is a limit of f with respect to the filter B and writing \(y = \lim_{B} f\) , we use the special notation

\[
y = \lim _ {t \in I} f (x),
\]

and we say that y is a limit of f at the point a, or that \(f(x)\) tends to y as x tends to a. Similarly, instead of saying that y is a cluster point of f with respect to B, we say that y is a cluster point of f at the point a.

A consideration of the definition of continuity (§2, no. 1, Definition 1) Proposition 7 of no. 3 shows that:

PROPOSITION 9. A mapping \(f\) of a topological space \(\mathbf{X}\) into a topological space \(\mathbf{Y}\) is continuous at a point \(a \in \mathbf{X}\) if and only if \(\lim f(x) = f(a)\) .

COROLLARY 1. Let X, Y be two topological spaces, f a mapping of X into Y which is continuous at a point \(a \in X\) ; then, for every filter base B on X which converges to a, the filter base \(f(\mathfrak{B})\) converges to \(f(a)\) . Conversely if, for every ultrafilter U on X which converges to a, the ultrafilter base \(f(\mathfrak{U})\) converges to \(f(a)\) , then f is continuous at a.

The first assertion is an immediate consequence of Proposition 9. To prove the second, suppose that f is not continuous at a; then there is a neighbourhood W of \(f(a)\) in Y such that \(\overline{f}^{1}(W)\) does not belong to the filter B of neighbourhoods of a in X. Hence (§ 6, no. 4, Proposition 7) there is an ultrafilter U, finer than B, which does not contain \(\overline{f}^{1}(W)\) and therefore contains its complement \(A = X - \overline{f}^{1}(W)\) (§ 6, no. 4, Proposition 5); since \(f(A) \cap W = \emptyset\) , \(f(\mathfrak{U})\) does not converge to \(f(a)\) .

COROLLARY 2. Let g be a mapping of a set Z into a topological space X, which has a limit a with respect to a filter F on Z; then if the map f: X → Y is continuous at a, the composition f\(\circ\)g has f(a) as a limit point with respect to F.

